\documentclass[11pt,a4paper]{article}

% --- Pacotes de Idioma e Codificação ---
\usepackage[utf8]{inputenc}
\usepackage[T1]{fontenc}
\usepackage[portuguese]{babel}

% --- Design e Layout ---
\usepackage[margin=2.5cm]{geometry}
\usepackage{charter} % Fonte profissional e legível
\usepackage{lastpage}
\usepackage{fancyhdr}
\usepackage[table]{xcolor}
\usepackage{tabularx}
\usepackage{booktabs}
\usepackage{xcolor}
\usepackage{enumitem}
\usepackage{amssymb}

% --- Configuração de Cores ---
\definecolor{primaryblue}{RGB}{0, 51, 102}

% --- Cabeçalho e Rodapé ---
\pagestyle{fancy}
\fancyhf{}
\fancyhead[L]{\small \textbf{Requisitos Consolidados} | EcoRide Solutions}
\fancyhead[R]{\small \thepage\ de \pageref{LastPage}}
\fancyfoot[L]{\small \textit{Confidencial - Apenas para uso interno}}
\renewcommand{\headrulewidth}{0.4pt}

% --- Início do Documento ---
\begin{document}
	
	% Título Principal
	{\noindent\LARGE\bfseries\color{primaryblue} Requisitos Consolidados}
	
	% ========== REQUISITOS GERAIS ==========
	\section*{Requisitos Gerais}
	
	\renewcommand{\arraystretch}{1.5}
	\begin{tabularx}{\textwidth}{@{}|l|X|@{}}
		\hline
		\rowcolor{gray!10} \textbf{Número} & \textbf{Descrição} \\ \hline
		\textbf{1} & O sistema deve calcular automaticamente o valor final da reparação com base na soma da reparações feitas, peças utilizadas e preços das reparações definidos pelo gerente. \\ \hline
		\textbf{2} & Sempre que uma OS for colocada nos estados: "Pendente de aprovação de orçamento", "Pendente Pagamento" ou "A aguardar peças", deverá ser gerado automaticamente um alerta (que contém a referência da OS e o estado) para a secretária, que deverá ficar guardado na sua lista de alertas. Cada alerta deverá ter um campo que indique se já foi tratado ou não. \\ \hline
		\textbf{3} & Deverá haver um registo de qual mecânico realizou o diagnóstico e qual realizou a reparação da OS. \\ \hline
		\textbf{4} & O sistema deve impedir que dois mecânicos iniciem a mesma OS simultaneamente, atribuindo-a a um mecânico de cada vez. Caso um mecânico selecione uma OS que já está a ser tratada, deverá ser impedido. \\ \hline
		\textbf{5} & Sempre que uma peça for colocada no estado "Possível defeito" ou tiver a sua quantidade mínima atingida, deverá ser gerado automaticamente um alerta (que contém o motivo do alerta e a peça em questão) para o gestor de stock, que deverá ficar guardado na sua lista de alertas. Cada alerta deverá ter um campo que indique se já foi tratado ou não. \\ \hline
	    \textbf{6} & Quando uma peça no estado "Possível Defeito" for colocada no estado "Pendente de Devolução", a quantidade em stock dessa peça deverá diminuir em 1 unidade. \\ \hline
 	    \textbf{7} & Quando uma peça no estado "Pendente de Devolução" for colocada no estado "Devolvida ao fornecedor", uma unidade dessa peça deverá ser adicionada ao stock sem qualquer custo adicional. \\ \hline
  	    \textbf{8} & Quando uma peça no estado "Pendente de Devolução" for colocada no estado "Inválida para Devolução", não deverá existir qualquer consequência no stock ou nos movimentos. \\ \hline
  	    \textbf{9} & O sistema deverá separar os utilizadores definindo diferentes perfis, consoante o seu cargo, disponibilizando as respetivas capacidades. \\ \hline
	\end{tabularx}
	
	\vspace{1cm}
	
	% ========== GERENTE ==========
	\section*{Requisitos do Gerente}
	
	\begin{tabularx}{\textwidth}{@{}|l|X|@{}}
		\hline
		\rowcolor{gray!10} \textbf{Número} & \textbf{Descrição} \\ \hline
		\textbf{1} & O gerente deverá ter acesso para fazer todas as ações que a secretária, gestor de stock e mecânicos podem realizar. \\ \hline
		\textbf{2} & O gerente deverá poder adicionar, remover e editar funcionários. \\ \hline
		\textbf{3} & O registo de funcionários deverá ser composto obrigatoriamente por: nome, morada (número de porta, rua, localidade e código-postal), telemóvel, email, data de nascimento, NISS, NIF, NUS, IBAN e salário (valor pago por hora e valor pago ao fim do mês, tanto líquido como bruto). \\ \hline
		\textbf{4} & O sistema deve permitir registar horas extraordinárias e calcular automaticamente o valor mensal a pagar de cada um dos funcionários (somando o valor pago pelas horas extraordinárias com o salário base). \\ \hline
		\textbf{5} & O sistema deverá permitir ao gerente consultar o estado financeiro da empresa, apresentando os movimentos monetários realizados (pagamentos de encomendas, pagamento de salários e receção de pagamentos). \\ \hline
		\textbf{6} & O sistema deverá ainda disponibilizar mecanismos de filtragem por intervalo de datas e por tipo de gasto ou receita, permitindo visualizar os valores totais por categoria, como salários, gastos com peças, lucros com mão de obra e lucros provenientes da venda de peças. \\ \hline
		\textbf{7} & O gerente deverá poder alterar a tabela de reparações, podendo adicionar, remover ou editar novas entradas à mesma. \\ \hline
		\textbf{8} & Cada entrada deverá ser composta por: uma referência única, uma nomenclatura, uma descrição detalhada e um preço de execução. \\ \hline
	\end{tabularx}
	
	\vspace{1cm}
	
	% ========== GESTOR DO STOCK ==========
	\section*{Requisitos do Gestor do Stock}
	
	\begin{tabularx}{\textwidth}{@{}|l|X|@{}}
		\hline
		\rowcolor{gray!10} \textbf{Número} & \textbf{Descrição} \\ \hline
		\textbf{1} & O gestor de armazém poderá consultar as peças existentes no sistema, as peças em stock, as peças com possível defeito em stock, as peças pendentes de devolução e as peças devolvidas ao fornecedor. Para isto, poderá filtrar por estado (se existe no sistema, stock, com possível defeito, pendente de devolução, devolvida ao fornecedor ou inválida para devolução), fornecedor, referência, data de receção (para as que estejam em stock), tipo ou preço. \\ \hline
		\textbf{2} & O gestor de armazém poderá adicionar, remover e editar o stock de peças existentes. Para além disso, também poderá adicionar, remover e editar fornecedores. \\ \hline
		\textbf{3} & O sistema deve permitir registar fornecedores obrigatoriamente com: nome e contacto (email ou número de telefone). \\ \hline
		\textbf{4} & O sistema deve permitir o registo de peças obrigatoriamente com: fornecedor, referência (do fornecedor), preço de compra, preço de venda, nível mínimo aceitável em stock, tempo de garantia e marca. \\ \hline
		\textbf{5} & O sistema deve permitir registar entradas de peças obrigatoriamente com: referência, quantidade, fornecedor, preço de compra e venda e data de receção. \\ \hline
		\textbf{6} & O sistema deve gerar listas de peças a encomendar com base em níveis mínimos, consumo e stock atual. \\ \hline
		\textbf{7} & O sistema deve permitir registar e editar peças devolvidas ao fornecedor, com data, motivo e estado da devolução ("Pendente de Devolução", "Devolvida ao Fornecedor" ou "Inválida para Devolução"). \\ \hline
	\end{tabularx}
	
	\vspace{1cm}
	
	% ========== ASSISTENTE ADMINISTRATIVA ==========
	\section*{Requisitos da Assistente Administrativa}
	
	\begin{tabularx}{\textwidth}{@{}|l|X|@{}}
		\hline
		\rowcolor{gray!10} \textbf{Número} & \textbf{Descrição} \\ \hline
		\textbf{1} & A secretária deverá poder consultar todas as OS consoante os seus estados ("Pendente Diagnóstico", "Pendente Reparação", "A aguardar peças", "Pendente aprovação do orçamento", "Pendente Pagamento", "Paga"), e deverão existir filtros que permitam filtrar as OS por estado, data, cliente e mecânico. \\ \hline
		\textbf{2} & A secretária deverá poder adicionar, remover e editar clientes, trotinetes e OS. \\ \hline
		\textbf{3} & O registo de clientes deverá ser composto obrigatoriamente por: nome, email, telemóvel e NIF. \\ \hline
		\textbf{4} & O registo de trotinetes deverá ser composto obrigatoriamente por: marca, modelo, número de série, tipo de motor e cliente. \\ \hline
		\textbf{5} & O sistema deverá registar a criação de uma OS obrigatoriamente com: os dados do cliente, os dados técnicos da trotinete, uma descrição do problema e, opcionalmente, por acessórios e estado à entrada (com notas textuais). Para todas as OS que forem criadas, deverá ser atribuído o estado de "Pendente Diagnóstico" e deverá ser criada uma referência única para a mesma.\\ \hline
		\textbf{6} & O sistema deve permitir registar a aprovação do orçamento por parte de um cliente somente em OS que estejam no estado "Pendente aprovação do orçamento", antes de permitir que o mecânico continue a reparação. Esta aprovação deverá colocar a OS num estado de "Pendente Reparação". \\ \hline
		\textbf{7} & Para uma OS no estado "Pendente Pagamento" \space deverá ser registado quando um cliente foi notificado sobre o término da reparação da trotinete. \\ \hline
		\textbf{8} & O sistema deve permitir registar a efetuação do pagamento e o respetivo método utilizado numa OS no estado "Pendente Pagamento" (numerário, multibanco, MBWay). Caso sejam ambos registados, a OS deverá ficar no estado "Paga". \\ \hline
	\end{tabularx}
	
	\vspace{1cm}
	
	% ========== MECÂNICO ==========
	\section*{Requisitos do Mecânico}
	
	\begin{tabularx}{\textwidth}{@{}|l|X|@{}}
		\hline
		\rowcolor{gray!10} \textbf{Número} & \textbf{Descrição} \\ \hline
		\textbf{1} & O mecânico poderá escolher de entre duas listas de OS: "Pendentes de diagnóstico" e "Pendentes de Reparação". \\ \hline
		\textbf{2} & Quando um mecânico selecionar uma OS "Pendente Diagnóstico", deverá ser possível escolher quais as reparações que ele acha necessárias realizar a partir da tabela de reparações que estará disponível, poderá consultar (procurando por referência ou fornecedor) e associar peças do stock, e também deverá poder consultar rapidamente o histórico de reparações da trotinete (mostrando toda a informação das OS que já foram levantadas sobre ela). \\ \hline
		\textbf{3} & Quando um mecânico selecionar uma OS "Pendente Diagnóstico" (e selecionar as várias reparações e peças usadas), deverá ser atribuído automaticamente o estado de "Pendente de aprovação de orçamento" e o mecânico não poderá proceder com a sua reparação.\\ \hline
		\textbf{4} & Quando um mecânico selecionar uma OS no estado "Pendente de reparação", deverá poder selecionar quais as reparações que ele executou a partir da tabela de reparações e poderá consultar (procurando por referência ou  fornecedor) e associar as peças que usou. Também deverá poder consultar rapidamente o histórico de reparações da trotinete (mostrando toda a informação das OS que já foram levantadas sobre ela). \\ \hline	
		\textbf{5} & Quando um mecânico selecionar uma OS no estado "Pendente de reparação" \space e uma ou mais peças da lista do diagnóstico não estiver disponível no stock, o sistema deverá disponibilizar a opção de requisição de encomenda das peças. A OS deverá ficar em estado de "A aguardar peças", e a reparação poderá prosseguir normalmente, mas a sua terminação deverá ser impedida até que o stock da peça necessário esteja disponível. Quando uma das peças estiver disponível no stock, a OS deverá ficar no estado "Pendente reparação". \\ \hline
		\textbf{6} & O sistema deve permitir ao mecânico adicionar novos problemas identificados durante a reparação, identificando-os com recurso à tabela de reparações. Esta ação deverá colocar o OS num estado de "Pendente aprovação do orçamento". \\ \hline
		\textbf{7} & Durante a realização de uma reparação, depois de adicionadas peças à lista de peças usadas na OS, o mecânico poderá reportar defeitos nas peças instaladas, atribuindo o estado "Possível defeito"\space à peça e removendo-a da lista de peças usadas, diminuindo o valor final da reparação. Caso não hajam mais peças iguais em stock, a OS deverá ficar em estado de "A aguardar peças", e a reparação poderá prosseguir normalmente, mas a sua terminação deverá ser impedida até que o stock da peça necessário esteja disponível. \\ \hline
		\textbf{8} & Quando um mecânico desejar terminar uma reparação, deverá verificar uma checklist obrigatória de segurança (luzes, pneus, aceleração, travagem, visor, teste de condução). \\ \hline
		\textbf{9} & O sistema deve permitir ao mecânico marcar a OS como concluída, ficando esta no estado "Pendente Pagamento". \\ \hline
	\end{tabularx}

	\section*{Requisitos da Investigação Operacional}
	
	\begin{tabularx}{\textwidth}{@{}|l|X|@{}}
		\hline
		\rowcolor{gray!10} \textbf{Número} & \textbf{Descrição} \\ \hline
		\textbf{1} & No sistema deverá ser possível nomear os itens que foram deixados na hora de entrega da trotinete, como por exemplo, carregador, cadeado e chave.  \\ \hline
		\textbf{2} & O sistema deverá ter um "Quadro de Estado em Tempo Real" \space acessível na receção. O mecânico deverá identificar a ordem de serviço associada e atualizar o seu estado para que a assistente tenha sempre a resposta imediatamente. \\ \hline
	\end{tabularx}

\end{document}
